4,674.51亿元在手订单是最强的收入可见性证据，也是最容易被误用的数字。它代表2025年末合并后上市公司的民船/海工手持合同，不代表中船集团整体，也不是完整军品订单；更不能用“除以三年”或直接折现的方式生成收入和利润。订单必须依次经过有效交付、会计确认、成本兑现、费用税负和少数股东归属。

\section{E05｜新签—在手—交付—收入—毛利漏斗}

2025年公司新签民船/海工237艘、3,050.67万DWT、1,758.36亿元；年末在手652艘、7,997.30万DWT、4,674.51亿元；全年交付161艘、1,367.75万DWT。2026计划交付168艘、1,650万DWT，并计划收入超过1,650亿元。\srcid{CF-03, CF-17}

\begin{exhibitbox}[订单到盈利的证据漏斗]
\centering
\begin{tikzpicture}[font=\small]
\foreach \w/\y/\c/\t in {
  13.5/0/navy/{在手订单：4,674.51亿元},
  11.5/-1.05/accentblue/{有效交付：2026 Base 1,620万DWT},
  9.5/-2.10/steelgrey/{Growth收入：2026 Base 1,552.41亿元},
  7.5/-3.15/riskamber/{合并毛利：2026 Base 296.42亿元},
  5.5/-4.20/riskgreen/{归母净利：2026 Base 172.76亿元}}
  {
    \fill[\c!16] (-\w/2,\y) -- (\w/2,\y) -- (\w/2-0.7,\y-0.78) -- (-\w/2+0.7,\y-0.78) -- cycle;
    \draw[\c,thick] (-\w/2,\y) -- (\w/2,\y) -- (\w/2-0.7,\y-0.78) -- (-\w/2+0.7,\y-0.78) -- cycle;
    \node at (0,\y-0.39) {\t};
  }
\end{tikzpicture}
\sourcenote{订单与交付CF-03/CF-17；2026 Base为冻结增长模型。漏斗不是法定履约率。}
\end{exhibitbox}

漏斗的收窄不是“订单流失”，而是不同量纲和会计层级。4,674.51亿元是多年度合同存量，2026收入只是当期履约确认；毛利再扣除费用、减值、税和少数股东后才成为归母。任何跳过中间层的订单估值都会高估确定性。

\section{收入桥：使用有效交付与收入强度代理，不伪造船型ASP}

模型把2025年船舶造修及海工收入1,312.79亿元除以交付1,367.75万DWT，得到95.9815亿元/百万DWT的“收入确认强度代理”。它同时包含造船进度、修船与海工收入，不是船东支付的每DWT船价，也不能跨船型解释价值。

\begin{equation*}
\text{Growth收入}_{t}=\text{有效交付DWT}_{t}\times95.9815\times\text{价格/结构指数}_{t}\times\text{确认系数}_{t}
\end{equation*}

基础情景2026采用1,620万DWT有效交付、1.04价格/结构指数与0.96确认系数，生成1,552.41亿元Growth收入，再加222.59亿元Base收入得到1,775亿元合并收入。这里的确认系数是相对2025收入强度的情景系数，不是公司披露的逐合同完工比例。

\begin{exhibitbox}[2026收入桥的披露与推导边界]
\small
\begin{tabularx}{\exhibitboxwidth}{L{3.4cm}R{2.4cm}L{4.0cm}X}
\toprule
\textbf{变量} & \textbf{Base} & \textbf{性质} & \textbf{验证证据} \\
\midrule
有效交付 & 1,620万DWT & 情景；低于1,650万DWT计划 & 正式交付、半年报收入 \\
收入强度代理 & 95.9815亿元/百万DWT & 2025披露数据推导；不是ASP & 2025核心收入与交付 \\
价格/结构指数 & 1.04 & 情景；反映船价与船型组合 & 船型合同额、收入强度 \\
确认系数 & 0.96 & 情景；不是法定履约比例 & 收入、合同资产、交付进度 \\
Growth收入 & 1,552.41亿元 & 上述四项机械结果 & 正式分产品收入 \\
Base收入 & 222.59亿元 & 分析拆分，不是披露分部 & 配套、机电及其他收入 \\
\bottomrule
\end{tabularx}
\sourcenote{CF-03、CF-17；冻结增长模型与分部预测桥。}
\end{exhibitbox}

这套桥的投资价值是可证伪：交付偏慢影响量，船型或签价变差影响价格/结构，合同资产堆积影响确认，正式收入可同时校验结果。它比“订单除以年数”更保守，也更适合在半年报后重跑。

\section{从收入到毛利：高船价仍要穿过65\%成本池}

公司称钢材与外购设备合计约占工业成本65\%，设备价格随订单景气上涨。高价订单的毛利还会受到人工、汇率、返工、质保与延期影响。2025年造船/修船/海工毛利率11.72\%，是2026改善的可比锚；说明会投资者提问中的17.48\%并非公司明确确认的核心业务披露，不被采用。\srcid{CF-03, CF-17}

\begin{exhibitbox}[订单价格到归母净利的驱动瀑布]
\small
\begin{tabularx}{\exhibitboxwidth}{L{3.1cm}L{4.2cm}L{3.5cm}X}
\toprule
\textbf{环节} & \textbf{正向驱动} & \textbf{负向驱动} & \textbf{财务落点} \\
\midrule
签价与产品结构 & 高船价、高端/绿色结构 & 低价旧单、结构下沉 & Growth收入与毛利 \\
生产效率 & 系列化、学习曲线、船位协同 & 返工、延误、熟练工瓶颈 & 毛利率与减值 \\
材料与设备 & 钢价稳定、集中采购 & 设备涨价、关键系统延迟 & 存货、成本、合同资产 \\
汇率与融资 & 出口收款与套保适配 & 人民币升值、客户融资失败 & 汇兑、收入确认与现金 \\
费用与归属 & 整合降费、少数股东占比改善 & 整合成本、关联交易与少数股东 & 归母净利润/EPS \\
\bottomrule
\end{tabularx}
\sourcenote{CF-03、CF-17；具体成本贡献为情景解释，不是逐合同披露。}
\end{exhibitbox}

所以，订单质量的最终验证项是毛利与现金，而不是金额本身。2026 Base Growth毛利率16.8\%显著高于2025年11.72\%，这正是模型最重要、也最需要正式中报验证的假设。

\section{中远项目：归属成立，但必须去重}

2025年报明确，中远集团约500亿元、87艘新造船项目均由600150下属相关子公司承建，上市公司聚合归属成立；同时项目含部分意向订单，具体确定/意向比例、船厂、船型、排期与利润率未披露。更重要的是，该项目已包含在公司聚合新签和在手披露中。\srcid{CF-03}

\begin{riskbox}[禁止重复叠加]
中远项目只作为客户关系与订单质量证据，所有情景的增量收入、EPS和目标价信用均为0。把约500亿元再次加到1,758.36亿元新签、4,674.51亿元在手或模型收入中，会造成实质性双重计算。
\end{riskbox}

同样的去重原则适用于江南180亿--190亿元双燃料集装箱船重大合同：它证明船厂能力与高端订单质量，但已属于聚合订单，不能另建合同价值加项。订单章节最终只保留一个闭环：聚合存量约束交付，交付与结构生成收入，毛利和现金决定估值信用。

\section{订单主张的失效条件}

若2026交付计划明显失速、合同资产和存货增加却没有相应收入、重大客户融资或撤单问题出现，订单可见性必须打折。若正式中报显示核心业务毛利不升反降，说明高船价可能被设备、人工、返工或低价旧单吞噬，不能继续以订单总额维持旧预测。

本章对行动的结论是：在手订单降低近期收入断崖风险，但没有提供现价之外的独立估值安全垫。新增资金必须等待漏斗中后段---毛利、归母与现金---出现可审计超预期，而不是在漏斗最上端重复计价。
